\begin{proof}[Proof of \cref{obj:lem:score-audit}]
\begin{enumerate}
\item Fix the schedule and design in the hypotheses. For \(i,j\in V\) and \(z\in\{0,1\}^{V}\), write
\[
N_i^-:=\{j\in V:(j,i)\in G\},
\qquad
X_j(z):=2z_j-1,
\qquad
X_j:=X_j(Z).
\]
The additive outcome map is
\[
Y_i(z)=a_i+t_i z_i+\sum_{j\in N_i^-}b_{ij}z_j.
\]
Index an audit mark by its recipient and source:
\[
W_{ij}:=W(j,i)\quad(j\ne i),
\qquad
H:=\{(j,i)\in G:W_{ij}=1\}.
\]
By \cref{obj:ass:assignment,obj:ass:audit,obj:ass:design-independence},
\[
0<q\le1,
\qquad
\operatorname{Law}(Z,W)
=
\left(\bigotimes_{j\in V}\operatorname{Bernoulli}(1/2)\right)
\otimes
\left(\bigotimes_{\substack{(j,i)\in V^2\\j\ne i}}
\operatorname{Bernoulli}(q)\right).
\]
We use \(\mathbb E_W\) for expectation over the audit marks with the assignment fixed, and \(\mathbb E_Z\) for expectation over the assignment. All the sums and expectations below are finite.

The evaluations of the scores in
\cref{obj:def:audit-score,obj:def:oracle-score} are
\[
\begin{aligned}
T_q(z,W)
&=\frac2n\sum_{i\in V}Y_i(z)
 \left(X_i(z)+q^{-1}\sum_{j\in N_i^-}W_{ij}X_j(z)\right),\\
T_G(z)
&=\frac2n\sum_{i\in V}Y_i(z)
 \left(X_i(z)+\sum_{j\in N_i^-}X_j(z)\right),
\end{aligned}
\qquad
T_q=T_q(Z,W),\quad T_G=T_G(Z).
\]
Since \(\mathbb E_W[W_{ij}]=q\), for every assignment \(z\),
\[
\mathbb E_W\!\left[\sum_{j\in N_i^-}W_{ij}X_j(z)\right]
=q\sum_{j\in N_i^-}X_j(z),
\qquad
\mathbb E_W[T_q(z,W)]=T_G(z).
\]

Under the half-Bernoulli assignment law, for \(j,k\in V\),
\[
\mathbb E_Z[X_j]=0,
\qquad
\mathbb E_Z[Z_kX_j]
=
\begin{cases}
\frac12,&k=j,\\
0,&k\ne j.
\end{cases}
\]
The first case follows from \(Z_j(2Z_j-1)=Z_j\); the second follows from independence and the zero mean of \(X_j\). Expanding the additive outcome therefore gives
\[
\mathbb E_Z[Y_i(Z)X_j]
=
\frac12\left(
t_i\mathbf1_{\{i=j\}}
+\sum_{k\in N_i^-}b_{ik}\mathbf1_{\{k=j\}}
\right).
\]
Because \(i\notin N_i^-\), summing this identity over the own coordinate and the incoming sources yields
\[
\mathbb E_Z\!\left[
Y_i(Z)\left(X_i+\sum_{j\in N_i^-}X_j\right)
\right]
=
\frac12\left(t_i+\sum_{j\in N_i^-}b_{ij}\right).
\]
Also, evaluating the outcome map at the two constant assignments gives
\[
Y_i(\mathbf1)-Y_i(\mathbf0)
=t_i+\sum_{j\in N_i^-}b_{ij}.
\]
Consequently,
\[
\mathbb E_Z[T_G]
=\frac1n\sum_{i\in V}\left(t_i+\sum_{j\in N_i^-}b_{ij}\right)
=\tau(\theta),
\qquad
\mathbb E_{P_{\theta,q}}[T_q]
=\mathbb E_Z\!\left[\mathbb E_W[T_q(Z,W)]\right]
=\tau(\theta).
\]
% lean: CausalSmith.Experimentation.ThinnedgraphAdditiveRiskFrontier.integral_auditScore_eq_tte

\item For a fixed assignment \(z\), subtracting the two score evaluations gives
\[
T_q(z,W)-T_G(z)
=
\frac{2}{nq}\sum_{i\in V}\sum_{j\in N_i^-}
Y_i(z)X_j(z)(W_{ij}-q).
\]
Each centered mark has mean zero and second moment
\[
\mathbb E_W[W_{ij}-q]=0,
\qquad
\mathbb E_W[(W_{ij}-q)^2]
=q(1-q)^2+(1-q)q^2
=q(1-q).
\]
For distinct off-diagonal pairs, independence gives
\[
\mathbb E_W[(W_{ij}-q)(W_{i'j'}-q)]=0
\qquad\bigl((j,i)\ne(j',i')\bigr).
\]
Thus variance is unchanged by subtracting the constant \(T_G(z)\), and the variance of the resulting weighted sum is
\[
\begin{aligned}
\operatorname{Var}_W(T_q(z,W))
&=\operatorname{Var}_W(T_q(z,W)-T_G(z))\\
&=\frac{4}{n^2q^2}\,q(1-q)
 \sum_{i\in V}\sum_{j\in N_i^-}Y_i(z)^2X_j(z)^2\\
&=\frac{4(1-q)}{n^2q}
 \sum_{i\in V}|N_i^-|\,Y_i(z)^2,
\end{aligned}
\]
where the last equality uses \(X_j(z)^2=1\). Combining this with the conditional mean already established gives the conditional second-moment identity
\[
\mathbb E_W[T_q(z,W)^2]
=
T_G(z)^2
+\frac{4(1-q)}{n^2q}
 \sum_{i\in V}|N_i^-|\,Y_i(z)^2.
\]
Empty incoming neighborhoods contribute zero. These identities also include \(q=1\), when every audit mark equals one almost surely.
% lean: CausalSmith.Experimentation.ThinnedgraphAdditiveRiskFrontier.integral_auditScore_sq_given_assignment

\item Integrating the preceding second-moment identity over the assignment gives
\[
\mathbb E_{P_{\theta,q}}[T_q^2]
=
\mathbb E_Z[T_G^2]
+\frac{4(1-q)}{n^2q}
 \sum_{i\in V}|N_i^-|\,\mathbb E_Z[Y_i(Z)^2].
\]
Both scores have mean \(\tau(\theta)\). Subtracting its square from the two second moments therefore yields
\[
\begin{aligned}
\operatorname{Var}_{P_{\theta,q}}(T_q)
&=\mathbb E_Z[T_G^2]-\tau(\theta)^2
 +\frac{4(1-q)}{n^2q}
  \sum_{i\in V}|N_i^-|\,\mathbb E_Z[Y_i(Z)^2]\\
&=\operatorname{Var}_Z(T_G)
 +\frac{4(1-q)}{n^2q}
  \sum_{i\in V}|N_i^-|\,\mathbb E_Z[Y_i(Z)^2].
\end{aligned}
\]
% lean: CausalSmith.Experimentation.ThinnedgraphAdditiveRiskFrontier.variance_auditScore_eq_oracle_add

\item To bound the oracle variance, augment each incoming neighborhood by its own coordinate. Define
\[
N_i:=\{i\}\cup N_i^-,
\qquad
\widetilde b_{ij}:=
\begin{cases}
t_i,&j=i,\\
b_{ij},&j\in N_i^-,
\end{cases}
\quad(i\in V,\ j\in N_i),
\qquad
\mu_i:=a_i+\frac12\sum_{j\in N_i}\widetilde b_{ij}.
\]
The two cases defining \(\widetilde b_{ij}\) are disjoint because the graph is off-diagonal. Substituting \(z_j=(1+X_j(z))/2\) into the outcome map gives
\[
Y_i(z)=\mu_i+\frac12\sum_{j\in N_i}\widetilde b_{ij}X_j(z),
\qquad
\tau(\theta)=\frac1n\sum_{i\in V}\sum_{j\in N_i}\widetilde b_{ij}.
\]
Separating the diagonal products in each row gives
\[
\begin{aligned}
\left(\sum_{j\in N_i}\widetilde b_{ij}X_j(z)\right)
\left(\sum_{k\in N_i}X_k(z)\right)
&=\sum_{j\in N_i}\widetilde b_{ij}\\
&\quad+\sum_{j\in N_i}\sum_{k\in N_i\setminus\{j\}}
\widetilde b_{ij}X_j(z)X_k(z),
\end{aligned}
\]
since \(X_j(z)^2=1\). Define, for every \(z\in\{0,1\}^{V}\),
\[
\begin{aligned}
T_{G,\mathrm{lin}}(z)
&:=\frac2n\sum_{i\in V}\mu_i\sum_{j\in N_i}X_j(z),\\
T_{G,\mathrm{quad}}(z)
&:=\frac1n\sum_{i\in V}\sum_{j\in N_i}
 \sum_{k\in N_i\setminus\{j\}}\widetilde b_{ij}X_j(z)X_k(z),
\end{aligned}
\qquad
\begin{aligned}
T_{G,\mathrm{lin}}&:=T_{G,\mathrm{lin}}(Z),\\
T_{G,\mathrm{quad}}&:=T_{G,\mathrm{quad}}(Z).
\end{aligned}
\]
Expanding the oracle score and subtracting its constant coefficient now gives
\[
T_G(z)-\tau(\theta)
=T_{G,\mathrm{lin}}(z)+T_{G,\mathrm{quad}}(z).
\]
% lean: CausalSmith.Experimentation.ThinnedgraphAdditiveRiskFrontier.oracleScore_centered_expansion

\item Since \(\mathbb E_Z[T_G]=\tau(\theta)\), the expansion gives
\[
\operatorname{Var}_Z(T_G)
=
\mathbb E_Z[T_{G,\mathrm{lin}}^2]
+\mathbb E_Z[2T_{G,\mathrm{lin}}T_{G,\mathrm{quad}}]
+\mathbb E_Z[T_{G,\mathrm{quad}}^2].
\]
Collecting the linear terms by their source coordinate gives
\[
T_{G,\mathrm{lin}}(z)
=
\frac2n\sum_{j\in V}
\left(\sum_{\substack{i\in V\\j\in N_i}}\mu_i\right)X_j(z).
\]

For any subsets \(E_1,E_2\subseteq V\), independence and centering imply
\[
\mathbb E_Z\!\left[
\left(\prod_{j\in E_1}(Z_j-\tfrac12)\right)
\left(\prod_{k\in E_2}(Z_k-\tfrac12)\right)
\right]
=
\begin{cases}
4^{-|E_1|},&E_1=E_2,\\
0,&E_1\ne E_2.
\end{cases}
\]
Indeed, when the sets agree, every factor is a square equal to \(1/4\). When they differ, a coordinate in their symmetric difference contributes an independent factor with mean zero. Using
\[
\prod_{j\in E_1}X_j
=2^{|E_1|}\prod_{j\in E_1}(Z_j-\tfrac12)
\]
and the corresponding identity for \(E_2\), we obtain
\[
\mathbb E_Z\!\left[
\left(\prod_{j\in E_1}X_j\right)
\left(\prod_{k\in E_2}X_k\right)
\right]
=\mathbf1_{\{E_1=E_2\}}.
\]
The empty product equals one, so these formulas include empty sets.

In particular,
\[
\mathbb E_Z[X_\ell X_jX_k]=0
\qquad(\ell,j,k\in V,\ j\ne k),
\]
because the singleton \(\{\ell\}\) differs from the two-element set \(\{j,k\}\). Expanding the collected linear term against the quadratic term expresses every summand of their cross moment as a constant times such a triple moment. Hence
\[
\mathbb E_Z[2T_{G,\mathrm{lin}}T_{G,\mathrm{quad}}]=0,
\qquad
\operatorname{Var}_Z(T_G)
=\mathbb E_Z[T_{G,\mathrm{lin}}^2]
+\mathbb E_Z[T_{G,\mathrm{quad}}^2].
\]
% lean: CausalSmith.Experimentation.ThinnedgraphAdditiveRiskFrontier.oracle_mixed_second_moment_zero

\item Set
\[
d_+:=d+1.
\]
The degree and coefficient-mass bounds in \cref{obj:def:model} imply
\[
|N_i|=1+|N_i^-|\le d_+,
\qquad
\left|\{i\in V:j\in N_i\}\right|
=1+\left|\{i\in V:(j,i)\in G\}\right|
\le d_+,
\]
and
\[
|a_i|+\sum_{j\in N_i}|\widetilde b_{ij}|
=
|a_i|+|t_i|+\sum_{j\in N_i^-}|b_{ij}|
\le1.
\]
In particular,
\[
\sum_{j\in N_i}|\widetilde b_{ij}|\le1,
\qquad
|\mu_i|
\le |a_i|+\frac12\left|\sum_{j\in N_i}\widetilde b_{ij}\right|
\le |a_i|+\sum_{j\in N_i}|\widetilde b_{ij}|
\le1.
\]

The sign orthogonality above, applied to singleton sets, gives
\[
\mathbb E_Z[T_{G,\mathrm{lin}}^2]
=
\frac4{n^2}\sum_{j\in V}
\left(\sum_{\substack{i\in V\\j\in N_i}}\mu_i\right)^2.
\]
For each \(j\), finite Cauchy--Schwarz and the coordinate multiplicity bound yield
\[
\left(\sum_{\substack{i\in V\\j\in N_i}}\mu_i\right)^2
\le
\left|\{i\in V:j\in N_i\}\right|
\sum_{\substack{i\in V\\j\in N_i}}\mu_i^2
\le
d_+\sum_{\substack{i\in V\\j\in N_i}}\mu_i^2.
\]
Summing and reversing the finite sums gives
\[
\begin{aligned}
\sum_{j\in V}
\left(\sum_{\substack{i\in V\\j\in N_i}}\mu_i\right)^2
&\le d_+\sum_{j\in V}\sum_{\substack{i\in V\\j\in N_i}}\mu_i^2\\
&=d_+\sum_{i\in V}|N_i|\mu_i^2\\
&\le d_+\sum_{i\in V}d_+
=nd_+^2.
\end{aligned}
\]
Since \(n\ge4\), scaling by \(4/n^2\) therefore gives
\[
\mathbb E_Z[T_{G,\mathrm{lin}}^2]\le\frac{4d_+^2}{n}.
\]
% lean: CausalSmith.Experimentation.ThinnedgraphAdditiveRiskFrontier.oracle_linear_scaled_second_moment_le

\item Define the collected ordered-pair coefficients for all \(j,k\in V\) by
\[
b^{\mathrm{pair}}_{jk}
:=
\sum_{\substack{i\in V\\j\in N_i,\ k\in N_i\setminus\{j\}}}
\widetilde b_{ij}.
\]
An empty sum equals zero; in particular,
\[
b^{\mathrm{pair}}_{jj}=0.
\]
Reversing the finite sums in the quadratic term gives
\[
nT_{G,\mathrm{quad}}
=\sum_{j\in V}\sum_{k\in V}b^{\mathrm{pair}}_{jk}X_jX_k.
\]
For \(j,k,j',k'\in V\) with \(j\ne k\) and \(j'\ne k'\), sign orthogonality gives
\[
\mathbb E_Z[X_jX_kX_{j'}X_{k'}]
=
\begin{cases}
1,&(j',k')=(j,k)\ \text{or}\ (j',k')=(k,j),\\
0,&\text{otherwise}.
\end{cases}
\]
Indeed, the two products correspond to the sets \(\{j,k\}\) and \(\{j',k'\}\), which agree precisely in those two orientations. Thus, for \(j\ne k\),
\[
\mathbb E_Z\!\left[
X_jX_k\sum_{j'\in V}\sum_{k'\in V}
b^{\mathrm{pair}}_{j'k'}X_{j'}X_{k'}
\right]
=b^{\mathrm{pair}}_{jk}+b^{\mathrm{pair}}_{kj}.
\]
Expanding the square and using \(b^{\mathrm{pair}}_{jj}=0\) for its diagonal terms consequently yields
\[
\mathbb E_Z[T_{G,\mathrm{quad}}^2]
=
\frac1{n^2}\sum_{j\in V}\sum_{k\in V}
b^{\mathrm{pair}}_{jk}
\left(b^{\mathrm{pair}}_{jk}+b^{\mathrm{pair}}_{kj}\right).
\]
% lean: CausalSmith.Experimentation.ThinnedgraphAdditiveRiskFrontier.oracle_quadratic_second_moment

\item The incidence condition defining \(b^{\mathrm{pair}}_{jk}\) is unchanged by swapping \(j\) and \(k\). Hence
\[
b^{\mathrm{pair}}_{jk}+b^{\mathrm{pair}}_{kj}
=
\sum_{\substack{i\in V\\j\in N_i,\ k\in N_i\setminus\{j\}}}
(\widetilde b_{ij}+\widetilde b_{ik}).
\]
The set of recipients in this sum is contained in
\(\{i\in V:j\in N_i\}\), whose size is at most \(d_+\). Cauchy--Schwarz therefore gives, for every \(j,k\in V\),
\[
\left(b^{\mathrm{pair}}_{jk}+b^{\mathrm{pair}}_{kj}\right)^2
\le
d_+\sum_{\substack{i\in V\\j\in N_i,\ k\in N_i\setminus\{j\}}}
(\widetilde b_{ij}+\widetilde b_{ik})^2.
\]
For \(j=k\), both sides are zero.

For \(j\in N_i\) and \(k\in N_i\setminus\{j\}\), the coefficient-mass bound gives
\[
|\widetilde b_{ij}|+|\widetilde b_{ik}|
\le\sum_{\ell\in N_i}|\widetilde b_{i\ell}|
\le1.
\]
The triangle inequality then gives the rowwise pair estimate
\[
(\widetilde b_{ij}+\widetilde b_{ik})^2
\le
\left(|\widetilde b_{ij}|+|\widetilde b_{ik}|\right)^2
\le
|\widetilde b_{ij}|+|\widetilde b_{ik}|.
\]
Summing this estimate and enlarging each inner sum by its nonnegative diagonal term yields
\[
\begin{aligned}
\sum_{j\in N_i}\sum_{k\in N_i\setminus\{j\}}
(\widetilde b_{ij}+\widetilde b_{ik})^2
&\le
\sum_{j\in N_i}\sum_{k\in N_i\setminus\{j\}}
\left(|\widetilde b_{ij}|+|\widetilde b_{ik}|\right)\\
&\le
\sum_{j\in N_i}\sum_{k\in N_i}
\left(|\widetilde b_{ij}|+|\widetilde b_{ik}|\right)\\
&=2|N_i|\sum_{j\in N_i}|\widetilde b_{ij}|\\
&\le2d_+.
\end{aligned}
\]
When \(N_i=\{i\}\), the sums over distinct coordinates are empty and the same bound applies.

Summing the pairwise Cauchy--Schwarz inequalities and reindexing by recipient now gives
\[
\begin{aligned}
\sum_{j\in V}\sum_{k\in V}
\left(b^{\mathrm{pair}}_{jk}+b^{\mathrm{pair}}_{kj}\right)^2
&\le
d_+\sum_{i\in V}\sum_{j\in N_i}
\sum_{k\in N_i\setminus\{j\}}
(\widetilde b_{ij}+\widetilde b_{ik})^2\\
&\le d_+\sum_{i\in V}2d_+
=2nd_+^2.
\end{aligned}
\]
Finally, expanding the square and exchanging \(j\) and \(k\) in one of the resulting sums gives
\[
\sum_{j\in V}\sum_{k\in V}
\left(b^{\mathrm{pair}}_{jk}+b^{\mathrm{pair}}_{kj}\right)^2
=
2\sum_{j\in V}\sum_{k\in V}
b^{\mathrm{pair}}_{jk}
\left(b^{\mathrm{pair}}_{jk}+b^{\mathrm{pair}}_{kj}\right).
\]
It follows that
\[
\sum_{j\in V}\sum_{k\in V}
b^{\mathrm{pair}}_{jk}
\left(b^{\mathrm{pair}}_{jk}+b^{\mathrm{pair}}_{kj}\right)
\le nd_+^2.
\]
% lean: CausalSmith.Experimentation.ThinnedgraphAdditiveRiskFrontier.oracle_quadratic_energy_le

\item Combining the exact quadratic second moment with this energy bound gives
\[
\mathbb E_Z[T_{G,\mathrm{quad}}^2]
\le\frac{nd_+^2}{n^2}
=\frac{d_+^2}{n}.
\]
Together with the vanishing cross moment and the linear bound, this proves
\[
\operatorname{Var}_Z(T_G)
\le\frac{4d_+^2}{n}+\frac{d_+^2}{n}
=\frac{5(d+1)^2}{n}.
\]
The unbiasedness and variance identity established above complete the proof.
% lean: CausalSmith.Experimentation.ThinnedgraphAdditiveRiskFrontier.oracle_variance_le
\end{enumerate}
% lean: CausalSmith.Experimentation.ThinnedgraphAdditiveRiskFrontier.score_audit
\end{proof}
